\section{A.2. Capping and its dependence on the other boundaries}
\begin{proposition}[Known with independent derivation: projective connection of a capped pant]\label{A2-schwarzian}
Let a hyperbolic pant have boundary lengths $\ell_0,\ell_1,\ell_\infty>0$, and canonically cap it at the three boundaries. Normalize the resulting sphere by placing the cap centres at $0,1,\infty$. Set
\[
a_j=\frac12\left(1+\frac{\ell_j^2}{4\pi^2}\right).
\]
The hyperbolic developing map extends as a complex projective structure to the punctured capped sphere, whose Schwarzian in the coordinate $z$ is
\[
T(z)=\frac{a_0}{z^2}+\frac{a_1}{(z-1)^2}
       +\frac{a_\infty-a_0-a_1}{z(z-1)}.
\]
If $u$ is any centred unit-disc coordinate on the cap at $0$, write $u(z)=c z(1+bz+O(z^2))$. Then
\[
b=\frac{a_0+a_1-a_\infty}{2a_0}.
\]
\end{proposition}
\begin{proof}
The projective connection is also the one in \cite[equation (1.3)]{AFirat2021}; we give the short derivation and the jet calculation needed here. Choose arclength $s$ on a boundary of length $\ell$, and a conformal collar coordinate $v=s+iy$ with boundary $y=0$. After an isometry of the upper half-plane, a hyperbolic developing map is $i\exp(v)$. The flat cylinder attached by arclength has punctured-disc coordinate $u=\exp(2\pi i v/\ell)$. Thus the continued projective developing map is $i u^{-i\ell/(2\pi)}$, with dilation monodromy. Its Schwarzian is $(1+\ell^2/(4\pi^2))/(2u^2)$. This gives a meromorphic projective connection with the stated double-pole coefficients at the three centres. On $\mathbf P^1$ a rational function with only these double poles and the required double-pole behavior at infinity is uniquely of the displayed form: a difference would be a holomorphic quadratic differential with at most simple poles at three points, whose degree is $-4+3<0$.

By the Schwarzian chain rule,
\[
T(z)=a_0\left(\frac{u'(z)}{u(z)}\right)^2+\{u,z\}.
\]
Here $u'/u=z^{-1}+b+O(z)$ and $\{u,z\}$ is holomorphic at zero. The simple-pole coefficient on the right is $2a_0b$, and that on the left is $a_0+a_1-a_\infty$. This proves the formula. Multiplication of $u$ by a rotation changes $c$ but not $b$.
\end{proof}

\begin{theorem}[Refuted: the independence asked in 67/068]\label{A2-independence}
In a canonically capped pant the cap at $0$ need not depend only on $\ell_0$. In particular, the pants with length triples $(2\pi,2\pi,2\pi)$ and $(2\pi,4\pi,2\pi)$ have different caps at $0$ in the common normalization $0,1,\infty$.
\end{theorem}
\begin{proof}
Hyperbolic pants exist and are unique for both positive length triples. Proposition \ref{A2-schwarzian} gives respectively $b=1/2$ and $b=5/4$. If the two pointed caps were the same domain, their centred disc coordinates would differ by a unit scalar, so their $b$ coefficients would agree. They do not. The corrected dependence is explicit in that proposition: at fixed $\ell_0$, the second jet already records $\ell_1^2-\ell_\infty^2$.
\end{proof}


% Secondary A.2(b) proof by the Part C agent.  Root may integrate this fragment.
% The external geometric inputs are infinite-grafting uniformization,
% the collar lemma, parameter dependence, and the weighted reduced-modulus theorem.
\subsection{The exact radius image and uniqueness of canonical caps}
\label{sec:A2b-secondary}

Let $Y$ be a connected compact marked Riemann surface of genus $g$, with
distinct marked points $x_1,\ldots,x_n$, where $2g-2+n>0$.
Choose a nonempty set $I$ of points that are to receive genuine caps;
the other marked points remain cusps. Fix a nonzero tangent coordinate
at each $x_i$, $i\in I$, so that conformal radii are positive numbers.
For $s=(s_i)_{i\in I}\in(0,\infty)^I$, put
\[
 \lambda_i=\sqrt{s_i},\qquad L_i=2\pi\lambda_i.
\]
Infinite-grafting uniformization gives the unique marked hyperbolic surface
with geodesic boundaries of these lengths whose infinite grafting has
conformal compactification $Y$. Denote its cap discs by $D_i(s)$, their
conformal radii by $r_i(s)$, and set $F_i(s)=\log r_i(s)$.
This use of the grafting theorem fixes the compact conformal surface,
as well as the indicated zero boundary lengths.

There is a related variational formula in the literature. In the sphere
setting treated by F\i rat, the on-shell action satisfies
$\partial S_{\mathrm{HJ}}/\partial\lambda_i=2\lambda_i\log r_i$
\cite[Appendix B, equation (B.1)]{AFirat2023}; after the change
$s_i=\lambda_i^2$, its derivative is $\log r_i$. We do not infer strict
convexity or the global range from that identity. The moving-seam energy
calculation below supplies these properties for every stable genus in
the present scope. We claim no priority for the variational viewpoint.

\begin{lemma}[Proved: regularity needed for variation]\label{A2b-regularity}
Locally in $(0,\infty)^I$, the canonical seams, their conformal cap maps,
and the hyperbolic and flat parts of the grafted metric depend
real-analytically on $s$. The logarithm of the metric density is $C^{1,1}$
across each seam and is smooth on either side. Its first parameter
variation is continuous across the seam and has the one-sided normal
derivatives used below.
\end{lemma}
\begin{proof}
The real-analytic inverse in the finite strata of the infinite-grafting
theorem gives a real-analytic family of marked hyperbolic surfaces, with
fixed cusp ends. Choose analytic Fenchel--Nielsen representatives locally.
Geodesic boundary collars and their normalized arc-length parameters
depend analytically on these representatives. Appending the standard
cylinders therefore gives an analytic family of complex charts, although
the metric, expressed in a chart crossing a seam, is only $C^{1,1}$.
The marked conformal isomorphisms of the compactifications with the fixed
$Y$ can be chosen real-analytically: apply parameter dependence for the
normalized Beltrami equation in these charts
\cite[\S5.7, Theorems 10--11, pp.~402--403]{AhlforsBers1960A2b}.
One can make the descent explicit on the universal cover of the marked
punctured surface. The normalized Beltrami solutions and the conjugated
deck transformations depend analytically on the parameters. Because
the marked conformal structure is fixed, these deck groups are conjugate
to the fixed group. Normalize that conjugacy by three distinct boundary
fixed points of deck transformations; its coefficients are analytic.
The resulting maps descend to the desired marked isomorphisms. In the
stable marked situation there is no nonidentity holomorphic automorphism
isotopic to the identity relative to the marked points. The resulting
maps and their parameter derivatives are smooth in the spatial variable
on a chart crossing the seam: its Beltrami coefficient is smooth, so
interior elliptic estimates applied to the Beltrami equation and its
parameter derivatives upgrade parameter dependence from the spaces in
that theorem to every local differentiability order. This proves the
claimed parameter regularity.

More explicitly, at a fixed seam use $\zeta=x+iy$, where $x$ has period
$2\pi\lambda_i$, $y>0$ is the flat side, and the initial metric is
\[
 g=\begin{cases}\sec^2y\,|d\zeta|^2,&y<0,\\
                  |d\zeta|^2,&y>0.
   \end{cases}
\]
The two logarithmic densities and their first derivatives agree at
$y=0$. Their second normal derivatives have finite one-sided values.
Differentiating their equality along the analytic moving seam gives
precisely the continuity and jump conditions calculated below; no
second derivative across the seam is assumed.
\end{proof}

\begin{theorem}[Proved: strict monotonicity of logarithmic cap radii]
\label{A2b-energy}
The Jacobian $dF/ds$ is symmetric positive definite everywhere.
Consequently $F$ is the gradient of a strictly convex real-analytic
function $S$ on $(0,\infty)^I$, and is injective. In particular an
admissible canonical cap system is determined by its conformal radii.
\end{theorem}
\begin{proof}
Take a tangent variation, and write
\[
 g_t=e^{2u_t}|dz|^2,\qquad h=\dot u_0,\qquad
 c_i=\dot\lambda_i/\lambda_i.
\]
The ratio defining $h$ is independent of the fixed local coordinate.
Let $X$ denote the hyperbolic core, including its cusp ends, and $C_i$
the flat cap cylinder. Differentiating the curvature equations gives
\[
 (\Delta_g-2)h=0\quad\hbox{on }X,\qquad
 \Delta_g h=0\quad\hbox{on }C_i.
\]
Use the collar coordinate of Lemma~\ref{A2b-regularity}, so the initial
cap coordinate is $w_i=\exp(i\zeta/\lambda_i)$.
Normalize the varying cap coordinate by its positive derivative in the
chosen tangent coordinate. The function
\[
 k_i=\operatorname{Re}\left.\frac{d}{dt}\right|_0\log w_{i,t},
 \qquad V_i=\lambda_i k_i
\]
is single-valued and harmonic on $C_i$. Since
$w_{i,t}(z)=z/r_i(t)+O(z^2)$ at its puncture,
\[
 V_i=-\lambda_i\dot F_i+O(e^{-y/\lambda_i}),\qquad
 h=c_i+O(e^{-y/\lambda_i})\quad(y\longrightarrow\infty).
\]
Differentiating
$g_t=\lambda_i(t)^2|dw_{i,t}/w_{i,t}|^2$ gives
\begin{equation}\label{A2b-cap-variation}
 h=c_i-\partial_y V_i.
\end{equation}
The equation $|w_{i,t}|=1$ shows that $V_i(x,0)$ is the normal velocity
of the seam in the direction from the core towards the flat cylinder.
The initial logarithmic density is $-\log\cos y$ on the core and $0$
on the flat side. Differentiating equality of its first derivatives
along the moving seam gives
\begin{equation}\label{A2b-jump}
 (\partial_y h)_+-(\partial_y h)_-=V_i.
\end{equation}
The sign follows also directly from the second derivatives $0$ and $1$
of those two logarithmic densities. Here $+$ denotes the flat side.

All the following integrals converge. The flat-end estimates above give
exponential decay of derivatives. At a fixed cusp, writing its canonical
coordinate as $q_t(z)=a_tz+O(z^2)$ in
$|dq_t/q_t|^2/(\log|q_t|)^2$ shows that
$h=O(1/|\log|z||)$ and that the truncated boundary terms tend to zero.
Green's identity, first on truncated pieces and then in the limit, yields
\begin{align*}
 E_h&:=\int_X\bigl(|\nabla h|^2+2h^2\bigr)dA
       +\sum_i\int_{C_i}|\nabla h|^2dA
       =-\sum_i\int_{\Gamma_i}hV_i\,dx,\\
 E_V&:=\sum_i\int_{C_i}|\nabla V_i|^2dA
       =\sum_i\int_{\Gamma_i}hV_i\,dx
          -\sum_i c_i\int_{\Gamma_i}V_i\,dx.
\end{align*}
In the first identity the core outward normal is $+\partial_y$ and
the cylinder outward normal is $-\partial_y$; use
\eqref{A2b-jump}. In the second use
\eqref{A2b-cap-variation}. The circular mean of a bounded harmonic
function on a half-cylinder is constant, so
\[
 \int_{\Gamma_i}V_i\,dx=-2\pi\lambda_i^2\dot F_i.
\]
Adding the identities proves the exact formula
\begin{equation}\label{A2b-positive-energy}
 \boxed{\quad
  \pi\sum_i\dot s_i\dot F_i
       =E_h+E_V.\quad}
\end{equation}
The right side is positive for a nonzero tangent vector: if it is zero,
$h=0$ on the nonempty hyperbolic core; continuity and the bounded
harmonic Dirichlet problem on each cylinder then give $h=0$ there,
and hence every $c_i=0$.

For two tangent variations $a,b$, apply the same integrations by parts
bilinearly. Their left side is
$\pi\sum_i\dot s_i^{,a}\dot F_i^{,b}$ and their right side is
\[
 \int_X(\langle\nabla h_a,\nabla h_b\rangle+2h_ah_b)dA
 +\sum_i\int_{C_i}
   (\langle\nabla h_a,\nabla h_b\rangle+
    \langle\nabla V_{i,a},\nabla V_{i,b}\rangle)dA,
\]
which is symmetric in $a,b$. This proves symmetry of $dF/ds$.
Since the positive orthant is simply connected, $F=dS$ for a function
$S$, and \eqref{A2b-positive-energy} makes $S$ strictly convex.
Finally, for distinct $s,t$, integrate along their straight segment:
\[
 (t-s)\cdot(F(t)-F(s))
 =\int_0^1 (t-s)^T(dF)_{s+v(t-s)}(t-s)\,dv>0.
\]
Thus $F$ is globally injective, as claimed.
\end{proof}

To specify the image, let $p\in[0,\infty)^I\setminus\{0\}$.
On $Y$ take the Jenkins--Strebel differential with cylinder circumferences
$2\pi\sqrt{p_i}$ at the points with $p_i>0$, all other marked points
being zero-weight marked points on its critical graph. Write $U_i(p)$
for its disc domains and $R_i(p)$ for their conformal radii. Define
\begin{equation}\label{A2b-support}
 K(p)=\sum_{p_i>0}p_i\log R_i(p),\qquad K(0)=0.
\end{equation}
Multiplying $p$ by a positive scalar rescales the differential and does
not change its disc domains. Equivalently, the weighted reduced-modulus
theorem gives
\begin{equation}\label{A2b-support-extremal}
 K(p)=\sup_{(D_i)}\sum_i p_i\log r_i(D_i),
\end{equation}
where the supremum ranges over pairwise disjoint pointed conformal
discs avoiding the other marked points; a zero-weight term is omitted.
Solynin's Theorem~2.1, equation (2.27), specializes to this assertion
\cite[\S1.2, (1.2), p.~9; \S2.3, (2.2), p.~23;
 Theorem 2.1, (2.27), p.~32]{Solynin1999A2b}:
take only the homotopy classes of circles about the positive-weight
points, no $G_j$ configurations in his equation (2.2), and use the
reduced disc modulus $(2\pi)^{-1}\log r_i$. His trajectory weights
$a_i=2\pi\sqrt{p_i}$ give the common factor $2\pi$ on both sides.
Zero-weight marked points are his allowed poles of order at most one;
such a pole may be removable when it meets a zero of the differential.
If one requires
nonzero competing discs at the zero-weight points too, the supremum
is unchanged: shrink the positive-weight extremal discs concentrically
by factors tending to $1$ and insert small discs at the remaining
points. This also
explains why the square roots of $p_i$, rather than $p_i$, are the
Strebel circumference weights.

The support value is computable directly from the unique differential,
without solving another hyperbolic capping problem. Write its principal
part at $x_i$ as $-p_i\,dz_i^2/z_i^2$. Choose a regular point $b_i$ of
the adjacent critical graph and the branch of
$\sqrt{-q_p/p_i}$ with residue $1$ at $x_i$. Then
\[
 \log R_i(p)=\lim_{z\to x_i}
 \left(\log|z_i(z)|-
  \operatorname{Re}\int_{b_i}^{z}\sqrt{-q_p/p_i}\right).
\]
The integral is taken inside the disc domain; its real part is
single-valued, and changing $b_i$ along the adjacent critical boundary
changes the integral only by an imaginary constant. This follows by
integrating $dw_i/w_i=\sqrt{-q_p/p_i}$ and normalizing $|w_i|=1$ on
that boundary.

\begin{lemma}[Proved: uniform cap bounds and recession]\label{A2b-recession}
There are constants $B_i<\infty$, depending only on the marked $Y$ and
the chosen tangent coordinates, such that every embedded pointed disc
avoiding the other marked points has conformal radius at most $B_i$.
For a canonical cap one has
\begin{equation}\label{A2b-collar-bound}
 r_i(s)\leq B_i\exp\left(
   -\frac{\arctan(\operatorname{csch}(\pi\sqrt{s_i}))}{\sqrt{s_i}}
                         \right).
\end{equation}
Moreover, for every $s_0>0$ and $p\geq0$,
\begin{equation}\label{A2b-recession-formula}
 \lim_{t\to\infty}\frac{S(s_0+tp)-S(s_0)}t=K(p).
\end{equation}
\end{lemma}
\begin{proof}
For the radius bound in genus zero, put the given point at $0$ and two
other marked points at $1,\infty$; Koebe's quarter theorem gives the
bound $4$, with the appropriate tangent-coordinate factor. In genus one,
lift the embedded disc to the universal covering flat plane. Its area
is at most the area $A$ of the torus, whereas the area formula for a
univalent map $f(z)=a_1z+\cdots$ gives area at least $\pi|a_1|^2$.
In higher genus, lift to the universal covering disc and apply
Schwarz's lemma. These arguments also cover extra marked points.

Double the hyperbolic surface across its geodesic boundaries. The collar
lemma gives an embedded one-sided collar of width
$w_i=\operatorname{arcsinh}(1/\sinh(L_i/2))$ on the original surface.
In the conformal coordinate above its width is
$a_i=\arctan(\sinh w_i)=\arctan(\operatorname{csch}(\pi\lambda_i))$.
The cap inverse map therefore extends univalently from the unit disc
to the disc of radius $e^{a_i/\lambda_i}$. The preceding global bound
applied to this enlarged disc proves \eqref{A2b-collar-bound}.
In particular $F_i(s)\to-\infty$ as $s_i\downarrow0$, uniformly in
the other parameters, and always $F_i(s)\leq\log B_i$.

The area comparison of Lemma~\ref{A2-strebel-deficit} gives the uniform
weighted deficit
\[
 0\leq K(s)-s\cdot F(s)\leq 2g-2+n.
\]
Indeed, every peripheral loop has grafted-metric length at least its
geodesic boundary length: project any flat-cylinder excursions to the
boundary and use the hyperbolic geodesic minimum. Cauchy--Schwarz on
each circle of the Strebel cylinders then compares their truncated
areas with the grafted area. The latter has hyperbolic part
$2\pi(2g-2+n)$ and the difference of its reduced flat areas is exactly
$2\pi s\cdot(F(s)-\log R(s))$.
Consequently, for $p>0$, the derivative
$\frac{d}{dt}S(tp)=p\cdot F(tp)$ tends to $K(p)$.
Ces\`aro averaging gives $S(tp)/t\to K(p)$.
This growth limit is unchanged by the fixed positive shift $s_0$:
the function $T(s)=S(s)-\sum_i(\log B_i)s_i$ is coordinatewise
nonincreasing, and, for $C$ such that $s_0\leq Cp$,
\[
 T((t+C)p)\leq T(s_0+tp)\leq T(tp).
\]
Dividing by $t$ proves \eqref{A2b-recession-formula} for $p>0$.
For a nonnegative $p$ with zero entries, feasibility of the cap discs
and \eqref{A2b-support-extremal} imply
$p\cdot F(s)\leq K(p)$, hence the upper bound $K(p)$ for the
recession limit. Set $b_i=\log B_i$. The coordinate derivative bounds
give, for $\epsilon>0$,
\[
 S(s_0+tp)\geq S(s_0+t(p+\epsilon\mathbf1))
                       -t\epsilon\sum_i b_i.
\]
The lower recession bound is therefore
$K(p+\epsilon\mathbf1)-\epsilon\sum_i b_i$.
From the supremum formula and the coordinate upper bounds,
$K(p+\epsilon\mathbf1)\to K(p)$: the upper bound is
$K(p)+\epsilon\sum b_i$, and any fixed nearly maximizing disc system
gives the required lower limit. Letting $\epsilon\downarrow0$ proves
the assertion, including the boundary directions.
\end{proof}

\begin{theorem}[Proved: exact image of the shadow map]\label{A2b-image}
The canonical cap radius map is a real-analytic diffeomorphism
\[
 (0,\infty)^I\longrightarrow\exp\mathcal L_Y,
 \qquad
 \mathcal L_Y=
 \left\{x\in\R^I:\ p\cdot x<K(p)
       \text{ for every }p\in[0,\infty)^I\setminus\{0\}\right\}.
\]
In particular its logarithmic image is open, convex, and downward closed.
It is precisely the set of conformal-radius vectors realized by systems
of pointed discs with pairwise disjoint analytic closures avoiding all
the other marked points. The boundary
support values are the weighted Strebel extrema \eqref{A2b-support}.
\end{theorem}
\begin{proof}
For a canonical system, one can enlarge any one cap a little into its
embedded hyperbolic collar while keeping every other cap fixed. For
any nonzero $p\geq0$, enlarge an index with $p_i>0$; the supremum formula
then gives the strict inequality $p\cdot F(s)<K(p)$.

Fix $x\in\mathcal L_Y$ and put $G(s)=S(s)-x\cdot s$.
Its sublevel sets are bounded. Indeed, otherwise choose $s_k$ with
$\|s_k\|\to\infty$ and $G(s_k)$ bounded above. From a fixed $s_0>0$,
pass to a subsequence such that
$p_k=(s_k-s_0)/\|s_k-s_0\|\to p\geq0$, $\|p\|=1$.
For any fixed $t>0$, convexity gives
\[
 G(s_0+tp_k)\leq
 (1-t/\|s_k-s_0\|)G(s_0)
 +(t/\|s_k-s_0\|)G(s_k).
\]
Taking the limit yields $G(s_0+tp)\leq G(s_0)$ for every $t$,
contradicting the positive recession slope $K(p)-p\cdot x$.

Extend $G$ to the closed orthant by lower semicontinuous closure.
A supporting affine function at any interior point bounds it below,
so this extension never takes the value $-\infty$. Bounded sublevel
sets imply that it attains its infimum. No minimum can lie at a finite
boundary point $b$. If the boundary value is infinite this is immediate.
Otherwise consider $g(t)=G(b+t\mathbf1)$ for $t>0$.
The convex closure has $g(t)\to G(b)$ as $t\downarrow0$; this follows
also by comparing with the segments from $b$ to an arbitrary interior
point. On the other hand,
\[
 g'(t)=\sum_i(F_i(b+t\mathbf1)-x_i)\longrightarrow-\infty:
\]
at least one $b_i$ is zero, its summand tends uniformly to $-\infty$
by \eqref{A2b-collar-bound}, and all other summands have finite uniform
upper bounds. Thus $g(t)<G(b)$ for small $t>0$, a contradiction.
The minimizer lies in the positive orthant and satisfies $F(s)=x$.
Strict convexity makes it unique. Positive definiteness of the Jacobian
then gives the real-analytic inverse.

The displayed half-space description proves convexity and downward
closure. Openness follows from the local inverse theorem (or from
continuity of $K$ on a normalized compact simplex). Every canonical
system has disjoint analytic closures. Conversely, a system with
disjoint analytic closures can be enlarged independently a little;
the same strict-inequality argument places its logarithmic radii in
$\mathcal L_Y$, so the result already proved supplies its unique
canonical representative with those radii. This last assertion
identifies radii, not the arbitrary discs with their canonical
representative.
\end{proof}

\paragraph{Dependency checkpoint.}
The moving-seam identity, its strictness, the collar bound, and the
convex minimization argument are proved here. The use of the
Strebel support values requires the weighted reduced-modulus theorem
in the indicated specialization. The recession calculation uses the
explicit area comparison of Lemma~\ref{A2-strebel-deficit}; it requires
no inference from convergence of abstract Teichm\"uller points to
convergence of embedded cap domains.

\subsection{A quantitative extremal estimate and all large-boundary limits}

We retain the notation $s_i=(L_i/2\pi)^2$, $F_i(s)=\log r_i(s)$,
and $K(p)$ of Section~\ref{sec:A2b-secondary}. All radii use fixed tangent
coordinates on the fixed marked compactification $Y$.

\begin{lemma}[Proved: a uniform weighted-radius deficit]\label{A2-strebel-deficit}
Put $c=2g-2+n>0$. For every $s\in(0,\infty)^I$ and every pairwise
disjoint system of pointed conformal discs $E_i$ avoiding the other
marked points,
\[
 \sum_{i\in I}s_i\bigl(\log r_i(E_i)-F_i(s)\bigr)\leq c.
\]
Consequently
\[
 0\leq K(s)-s\cdot F(s)\leq c.
\]
The constants do not depend on the boundary lengths. Fixed cusp ends
are allowed.
\end{lemma}
\begin{proof}
Let $g_s$ be the metric which is hyperbolic on the bordered core and
flat on each appended cylinder. Its core has area $2\pi c$ by
Gauss--Bonnet, including the finite-area cusp ends. Every closed curve
in the peripheral free homotopy class of $x_i$ has $g_s$-length at least
$L_i$. Indeed, the map collapsing every flat cylinder orthogonally onto
its boundary is length nonincreasing. The resulting curve in the
hyperbolic core has length at least its geodesic representative, of
length $L_i$. This argument also covers a curve originally contained
entirely in the cylinder.

Let $\zeta_i:E_i\to\D$ be its centred uniformization. On the annulus
$\epsilon<|\zeta_i|<R$, $R<1$, write the metric as
$\rho_i^2|d\zeta_i|^2$. Each circle is peripheral, so
\[
 \int_0^{2\pi}\rho_i(re^{i\theta})r\,d\theta\geq L_i.
\]
Cauchy--Schwarz, followed by integration in $dr/r$, gives
\[
 \operatorname{Area}_{g_s}\{\epsilon<|\zeta_i|<R\}
 \geq\frac{L_i^2}{2\pi}\log(R/\epsilon).
\]
The test annuli are disjoint. Near $x_i$, the canonical cap coordinate
$w_i$ and the test coordinate satisfy
$w_i=(r_i(E_i)/r_i(s))\zeta_i+O(\zeta_i^2)$, up to a rotation.
The flat metric is $s_i|dw_i/w_i|^2$. Thus, after deleting the small
discs $|\zeta_i|\leq\epsilon$ at all active ends, the area of the whole
surface is
\[
 2\pi c+2\pi\sum_i s_i
 \bigl(F_i(s)-\log r_i(E_i)-\log\epsilon\bigr)+o(1).
\]
The error tends to zero for fixed $s$: the analytic coordinate change
has a nonzero first derivative, and its higher terms contribute
$O(\epsilon)$ to the boundary height. Compare this area with the sum
of the lower bounds, first let $\epsilon\downarrow0$, then let
$R\uparrow1$, and cancel the divergent terms. The first assertion
follows. Taking the supremum over the test systems gives the upper
bound in the second assertion by the weighted reduced-modulus
characterization of $K$. The lower bound holds because the canonical
caps themselves are a competing disc system.
\end{proof}

\begin{theorem}[Proved with the stated fixed-compactification hypothesis:
the cap limit]\label{A2-cap-limit}
Fix $Y$ and its marked points as above. Let $s^{(m)}>0$ and $t_m\to\infty$
satisfy $s^{(m)}/t_m\to p\in[0,\infty)^I\setminus\{0\}$.
For every $i$ with $p_i>0$, the pointed caps $D_i(s^{(m)})$ converge in
the Carath\'eodory kernel sense to the Jenkins--Strebel disc $U_i(p)$
with circumference weights $2\pi\sqrt{p_i}$. Their normalized conformal
maps converge locally uniformly, and
\[
 r_i(s^{(m)})\longrightarrow R_i(p).
\]
For $p_i=0$, the radii tend to zero. In particular the limit depends
on the relative growth of the lengths and on $Y$; it is not universally
a Koebe slit domain. No assertion of a fixed limit is made when the
marked compactification itself is allowed to diverge.
\end{theorem}
\begin{proof}
Every competing pointed disc has a finite uniform radius bound $B_i$.
For a sphere this is Lemma~\ref{A1-cap-bound} after normalizing two
other marked points to $1,\infty$. For a torus, lift to its flat
universal covering: injectivity downstairs bounds the area of the
lifted image by the torus area, and the area formula bounds its first
Taylor coefficient. For genus at least two, lift to the uniformizing
disc and use Schwarz's lemma. These arguments also give compactness
of normalized univalent maps whenever their radii stay bounded away
from zero. In the spherical case use the usual distortion bound for
normalized univalent functions; the bounded first coefficient makes
the family locally bounded. In the torus case the same distortion
bound applies to the lifts. In higher genus use normality of maps to
the uniformizing disc. A nonconstant limit remains injective, and
Hurwitz's theorem preserves omission of the other marked points.

Let $J=\{i:p_i>0\}$. Choose any fixed competing system $E_i$ for all
active points. Divide Lemma~\ref{A2-strebel-deficit} by $t_m$. Terms
with $i\notin J$ on the left side $\sum s_i^{(m)}F_i(s^{(m)})/t_m$
are bounded above by $\sum_{i\notin J}(s_i^{(m)}/t_m)\log B_i=o(1)$.
It follows that
\[
 \liminf_m\sum_{i\in J}\frac{s_i^{(m)}}{t_m}F_i(s^{(m)})
 \geq\sum_{i\in J}p_i\log r_i(E_i).
\]
The same estimate and the individual upper bounds imply a uniform
lower bound for each $F_i(s^{(m)})$, $i\in J$. Hence we may replace
$s_i^{(m)}/t_m$ by $p_i$ in this inequality. The supremum of the right
side is $K(p)$: approximate the positive-weight Strebel discs by compact
concentric subdiscs, and put sufficiently small discs at the remaining
points. Conversely, feasibility gives
$\sum_{i\in J}p_iF_i(s^{(m)})\leq K(p)$. Therefore this weighted sum
converges to $K(p)$.

From any subsequence extract convergent normalized cap maps at every
positive-weight point. Their nonconstant limits have disjoint images:
an overlap would put a common small open set into the images for all
large $m$, by local uniform convergence and the inverse mapping theorem.
Their radii are the limits of the first derivatives. Thus their disc
system attains the weighted reduced-modulus maximum $K(p)$.
The uniqueness clause in the Jenkins--Strebel extremal theorem makes
it precisely $\{U_i(p):i\in J\}$. This identifies every convergent
subsequence and proves the asserted convergence for the full sequence.
The locally uniform convergence of normalized univalent maps is
equivalent to the stated pointed kernel convergence.

Finally these positive-weight Strebel discs fill $Y$ outside their
critical graph, which has empty interior. If a zero-weight cap had
radius bounded away from zero along a subsequence, compactness would
give an additional nonempty open disc disjoint from all the limiting
positive-weight discs. Such an open set cannot lie in the critical
graph. This contradiction proves its radius tends to zero.
\end{proof}

\begin{corollary}[Known form, with explicit derivation: the Koebe and
orange limits]\label{A2-explicit-limits}
Normalize a capped pant by placing its centres at $0,1,\infty$.
If $L_0\to\infty$ and $L_1/L_0,L_\infty/L_0\to0$, then the cap at
$0$ tends to $\C\setminus[1,\infty)$, its conformal radius tends to
$4$, and its normalized conformal map tends to
\[
 k(u)=\frac{4u}{(1+u)^2},\qquad |u|<1.
\]
If instead $L_0/L_1,L_0/L_\infty\to1$ and the lengths tend to infinity,
the cap at $0$ tends to
\[
 Q_0^\circ=\{|z|<1,\ \operatorname{Re}z<1/2\},
 \qquad r(Q_0,0)=\frac4{3\sqrt3}.
\]
The other two limits follow by the permutations of $0,1,\infty$.
This is the precise limiting sense in which the orange slices are valid.
\end{corollary}
\begin{proof}
Up to the common positive scale, the three-pole Strebel quadratic
differential with squared circumference parameters $p_0,p_1,p_\infty$
has the rational expression
\[
 q_p=-\frac{p_\infty z^2+(p_1-p_\infty-p_0)z+p_0}
              {z^2(z-1)^2}\,dz^2.
\]
There is no accessory parameter: the difference between two such
quadratic differentials would have at most three simple poles on the
sphere and must vanish. For $p=(1,0,0)$ this is
$-dz^2/(z^2(1-z))$. Direct substitution of $z=k(u)$ gives
$k^*q_p=-du^2/u^2$. The map $k$ is conformal onto the stated slit
domain, as follows either by its explicit inverse using
$\sqrt{1-z}$ with value $1$ at $0$, or by its boundary image and
nonvanishing derivative in $\D$. Its derivative at zero is $4$.

For $p=(1,1,1)$ the numerator is $z^2-z+1$, with simple zeros
$a=e^{i\pi/3}$ and $b=e^{-i\pi/3}$. On the long unit-circle arc
joining $a$ to $b$, and on the vertical segment joining them,
$q_p$ is positive on the tangent directions. Together with the long
arc of $|z-1|=1$ joining the same vertices, these three symmetric
arcs form the critical graph. The complementary disc about $0$ is
$Q_0^\circ$. To check its radius explicitly, put
$m(z)=(z-a)/(z-b)$; it maps this lune to a sector of angle $2\pi/3$.
The branch $v(z)=(e^{-i\pi/3}m(z))^{3/2}$ maps that sector to the
upper half-plane and takes $0$ to $i$. Therefore
$U(z)=(v(z)-i)/(v(z)+i)$ is a centred disc uniformization.
Since $m'(0)/m(0)=i\sqrt3$, we obtain
$|U'(0)|=3\sqrt3/4$, giving the displayed radius. This also verifies
the cylinder structure of these complementary domains directly.
Theorem~\ref{A2-cap-limit} now proves both limits.
\end{proof}

\begin{corollary}[Known inequality with derivation: the image is not a box]
\label{A2-nonbox}
For the pants normalization $0,1,\infty$, every genuine admissible cap
system satisfies $r_0r_1<1$. Nevertheless the individual supremum of
each of $r_0,r_1,r_\infty$ over such systems is $4$. The set of
admissible radii is therefore not a Cartesian product of intervals.
\end{corollary}
\begin{proof}
For $p=(1,1,0)$ the quadratic differential is
$-dz^2/[z^2(z-1)^2]$. Its two disc domains are the half-planes separated
by $\operatorname{Re}z=1/2$; each has conformal radius $1$ at its marked
centre. The weighted reduced-modulus inequality gives $r_0r_1\leq1$,
the two-domain bound stated in the opened original
\cite[Part I, \S2, Theorem3, pp.~169--173]{ALavrentiev1934}:
his maps go from the domains to the unit disc, so their derivative
product bound is the reciprocal conformal-radius bound. It is strict for genuine
canonical caps: their disjoint analytic boundaries admit simultaneous
small enlargements into the nonempty hyperbolic core, which strictly
increase their radii. The individual suprema are at most $4$ by Koebe
and are approached by the one-dominant-boundary limits just proved.
A Cartesian product with these individual suprema would contain two
radii both larger than $1$, contradicting the product inequality.
\end{proof}

\subsection{Cusp norms and the collar question}
\begin{proposition}[Known with derivation: cusp norms and decorations]\label{A2-cusp-norm}
For a compact pointed complex curve $(Y,p_1,\ldots,p_n)$ with $2g-2+n>0$, the complete hyperbolic metric on $Y\setminus\{p_i\}$ defines canonical Hermitian norms on $T_{p_i}Y$. If $q_i$ is its cusp coordinate, then
\[
 ds^2=\frac{|dq_i|^2}{|q_i|^2\log^2|q_i|},\qquad
 \|v\|_{p_i}=|dq_i(v)|.
\]
The coordinate is unique up to a constant of modulus one. A horocycle of circumference $h_i$ has coordinate radius $\exp(-2\pi/h_i)$ wherever embedded. Penner's unrestricted decorations, understood as horocycles in the cusp cover and allowing their immersed images, form a principal $(\R_{>0})^n$-bundle. The map which assigns to $X$ the capped curve and decorations of circumferences $\ell_i(X)$ identifies its positive-length bundle with the one in Theorem \ref{A4-grafting}.
\end{proposition}
\begin{proof}
Normalize a primitive parabolic stabilizer of a cusp in the upper half-plane to $\zeta\mapsto\zeta+1$. Then $q=\exp(2\pi i\zeta)$ gives the displayed metric. Any two such normalized isometries differ by $\zeta\mapsto\zeta+a$ with $a\in\R$, giving precisely the unit scalar ambiguity and hence the intrinsic norm. The horocycle $\operatorname{Im}\zeta=y$ has length $1/y$; its $q$-radius is $e^{-2\pi y}$. Circumferences give independent positive real coordinates on the decoration fibre. Now apply Theorem \ref{A4-grafting}. This is Wolpert, \emph{Cusps and the family hyperbolic metric}, Section 1, Definition 1 \cite{AWolpertCusps2007}, together with the definition of Penner's decorated space \cite[Section 3, pp.~309--310]{APenner1987}. It applies in genus zero as well; the absence of a canonical smooth metric on the whole sphere causes no obstruction to a canonical cusp tangent norm. The norm does not assert that a unit tangent disc is realized by an embedded cap.
\end{proof}

\begin{theorem}[Proved: metric comparison after removing normal collars]\label{A2-collar}
Let $(M,g)$ be a compact hyperbolic surface with geodesic boundary and negative Euler characteristic. Remove pairwise disjoint normal collars of widths $d_i>0$ from selected boundary components, without changing the topological type, and call the remaining surface $M'$. Let $g'$ be its conformally compatible hyperbolic metric with geodesic boundary. Then $g'=e^{2u}g$ on $M'$ with $u>0$. More precisely $u$ is the unique smooth solution of
\[
\Delta_g u=e^{2u}-1\quad\hbox{on }M',\qquad
\partial_{n}u=\begin{cases}\tanh d_i&\hbox{on a new boundary},\\0&\hbox{on a retained boundary},\end{cases}
\]
where $\Delta=\operatorname{div}\operatorname{grad}$ and $n$ is the outward unit normal of $M'$. The new boundary corresponding to a former boundary of length $L_i$ has length strictly greater than $L_i\cosh d_i$ in $g'$.
\end{theorem}
\begin{proof}
Uniformization supplies $g'$ and its smooth extension to the analytic boundary. The curvature transformation gives $\Delta_g u=e^{2u}-1$. In normal coordinates at an old boundary, $g=dr^2+\cosh^2(r)ds^2$, with $r$ increasing into $M$. At $r=d_i$ the outward normal of the retained surface is $-\partial_r$, and the curvature vector of its unit tangent is $-\tanh(d_i)\partial_r$. Its outward geodesic curvature is therefore $\tanh d_i$. Under conformal scaling, this signed curvature becomes $e^{-u}(k_g-\partial_nu)$, giving the boundary condition.

If $u$ had a negative minimum in the interior, the equation would give $\Delta_g u<0$ there, a contradiction. If a negative minimum occurred on the boundary, write the equation as $\Delta_g u-c(x)u=0$ with $c(x)=(e^{2u}-1)/u>0$ (continuously set $c=2$ at zero). The boundary point lemma gives $\partial_nu<0$ for a nonconstant solution at such a minimum, contradicting the nonnegative prescribed derivative. A negative constant cannot solve the equation. Thus $u\geq0$. The strong maximum principle and the same boundary lemma imply $u>0$ everywhere, since a positive collar width excludes the identically zero solution.

For uniqueness, if $u_1,u_2$ solve the same problem, their difference $v$ satisfies $\Delta_gv=c(x)v$ with $c(x)>0$ and homogeneous Neumann boundary data. Integration by parts gives $-\int|\nabla v|^2=\int c v^2$, so $v=0$. Finally the old metric length of $r=d_i$ is $L_i\cosh d_i$ and the strict conformal factor inequality increases it. Thus the question of 67/021 has an intrinsic boundary-value answer; the old metric is not simply its restriction after the collar is removed.
\end{proof}


\subsection{Arbitrary smooth collars: an exact equation and the shape dependence}
The question on 67/021 says ``collar'' without prescribing its shape. We therefore treat any smooth embedded collar whose removal preserves the topological type, not only a collar of constant normal width. Curvature below is signed against the outward normal of the retained surface.

\begin{theorem}[Proved: the metric relation for every smooth collar]\label{A2-collar-general}
Let $(M,g)$ be a compact connected hyperbolic surface with geodesic boundary and $\chi(M)<0$. Let $M'\subset M$ be obtained by removing finitely many disjoint smooth open boundary collars, with $M'$ diffeomorphic to $M$. Write $k_g=g(\nabla_TT,n)$ on $\partial M'$, where $T$ is a unit tangent and $n$ is the outward unit normal. The unique conformally compatible hyperbolic metric with geodesic boundary is $g'=e^{2u}g$, where $u\in C^\infty(M')$ is the unique solution
\[
 \Delta_g u=e^{2u}-1,\qquad \partial_nu=k_g.
\]
Its area and boundary lengths are exactly
\[
 \int_{M'}e^{2u}\,dA_g=-2\pi\chi(M),\qquad
 L_{g'}(B)=\int_B e^u\,ds_g\quad(B\subset\partial M').
\]
If $k_g\ge0$ everywhere and is positive somewhere, then $u>0$ everywhere. In general the sign of $u$ is not determined by the fact that a collar was removed.
\end{theorem}
\begin{proof}
Uniformization of the bordered conformal surface supplies $g'$, smooth to the smooth boundary. The curvature formulas are
\[
 K_{e^{2u}g}=e^{-2u}(-1-\Delta_g u),\qquad
 k_{e^{2u}g}=e^{-u}(k_g-\partial_nu).
\]
Thus its defining conditions give the stated boundary-value problem. Conversely these formulas show that every smooth solution supplies the required metric. If $u_1,u_2$ are solutions, their difference $v$ satisfies $\Delta_gv=c(x)v$, $\partial_nv=0$, with
$c(x)=2\int_0^1e^{2((1-t)u_2+tu_1)}dt>0$. Multiplying by $v$ and integrating gives
$-\int|\nabla v|^2=\int cv^2$, so $v=0$.
Gauss--Bonnet for $g'$ gives its area, and conformal scaling gives the boundary-length formula. The maximum-principle and boundary-point argument in Theorem~\ref{A2-collar} uses only $k_g\ge0$, and proves the asserted strict positivity. The final assertion is proved by the following explicit family; no positivity is assumed when applying the displayed equation to a general collar.
\end{proof}

\begin{proposition}[Proved: first variation and a family refuting general monotonicity]\label{A2-collar-shape}
Let the boundary components of $(M,g)$ have arclength coordinates $s$, and choose smooth functions $a_i(s)>0$. In inward Fermi coordinates remove the collars
$0\le r<\varepsilon a_i(s)$, for sufficiently small $\varepsilon>0$, and let $u_\varepsilon$ be as in Theorem~\ref{A2-collar-general}. After any smooth identification with $M$ which is the identity at $\varepsilon=0$,
\[
 u_\varepsilon=\varepsilon h+O(\varepsilon^2)\quad\hbox{in }C^2,
 \qquad
 (\Delta_g-2)h=0,\quad \partial_nh=a_i-a_i''.
\]
Here the prime is differentiation in boundary arclength. In particular, take the unique hyperbolic pant all of whose boundary lengths are $2\pi$. On one boundary set
\[
 a_1(s)=1+\tfrac12\cos(ns),
\]
and on the other two set $a_2=a_3=1$. For every sufficiently large integer $n$, and then every sufficiently small $\varepsilon>0$, the resulting $u_\varepsilon$ takes a negative value at a point of the new first boundary. Thus even smooth analytic, everywhere inward collar removal need not increase the hyperbolic metric pointwise. The same metric takes a positive value somewhere.
\end{proposition}
\begin{proof}
In the original metric the collar is
$dr^2+\cosh^2r\,ds^2$. For the graph $r=\varepsilon a(s)$, its tangent acceleration in the $r$ direction is
$\varepsilon(a''-a)+O(\varepsilon^2)$, and its outward normal is $-\partial_r+O(\varepsilon)$. Hence
$k_g=\varepsilon(a-a'')+O(\varepsilon^2)$.
Pull the equation back to a fixed domain. Its derivative in $u$ at $(\varepsilon,u)=(0,0)$ is the Neumann operator
$h\mapsto((\Delta_g-2)h,\partial_nh)$. This is an isomorphism between the usual H\"older spaces: the form
$\int(|\nabla h|^2+2h^2)$ is coercive, so the weak solution exists uniquely, and local elliptic regularity up to a smooth boundary gives the asserted regularity. The implicit-function theorem gives smooth dependence on $\varepsilon$. Terms differentiating the pulled-back operator vanish on $u_0=0$. This proves the variation formula. It is independent of the chosen identification for the same reason.

We give a direct estimate establishing the sign in the displayed family. Let $h_0$ be the fixed solution with normal derivative $1$ on each boundary. Let $v_n$ solve $(\Delta_g-2)v_n=0$, with normal derivative $(1+n^2)\cos(ns)$ on the first boundary and zero on the other two. Then $h=h_0+v_n/2$.
Choose a fixed embedded collar $0\le r\le\delta$ of the first boundary and a smooth cutoff $\eta(r)$ equal to $1$ near zero and to $0$ near $\delta$. Define, extending by zero outside the collar,
\[
 w_n(r,s)=(n+n^{-1})\eta(r)e^{-nr}\cos(ns).
\]
It has exactly the same normal derivative as $v_n$. Put $R_n=(\Delta_g-2)w_n$. Since
\[
 \Delta_g=\partial_r^2+\tanh r\,\partial_r
                      +\operatorname{sech}^2r\,\partial_s^2,
\]
where $\eta=1$ the residual is
\[
 R_n=(n+n^{-1})e^{-nr}
       (n^2\tanh^2r-n\tanh r-2)\cos(ns).
\]
The terms involving derivatives of $\eta$ are supported away from zero and are exponentially small times a fixed polynomial in $n$. Using $|\tanh r|\le r$, changing variable $t=nr$, and using the bounded area-density factor $\cosh r$ on this collar, we obtain
\[
 \|R_n\|_{L^2(M)}^2=O(n),
 \qquad
 \int_0^\infty(t^2+t+2)^2e^{-2t}\,dt=\frac{23}{4}.
\]
All constants here depend only on the fixed pant and cutoff, not on $n$.

For clarity we justify the estimate on the correction, including its boundary. Double $M$ across its geodesic boundary. Its metric extends smoothly because $\cosh r$ is even. Extend $R_n$ evenly as an $L^2$ function on that closed hyperbolic surface. The weak solution of
$(\Delta-2)e_n=-R_n$ on the double is in $H^2$ and is even, so its restriction has homogeneous Neumann data and equals $v_n-w_n$. Integration by parts gives $\|e_n\|_2\le\|R_n\|_2/2$ and a bound for $\|\nabla e_n\|_2$ by a fixed multiple of $\|R_n\|_2$. The integrated Bochner identity, first for smooth functions and then by smooth approximation in $H^2$, on a curvature-minus-one surface is
\[
 \|\nabla^2e_n\|_2^2=\|\Delta e_n\|_2^2+\|\nabla e_n\|_2^2.
\]
Together with $\Delta e_n=2e_n-R_n$ it bounds the full $H^2$ norm by $C\|R_n\|_2$. Finally $H^2$ embeds in continuous functions on a compact surface: in each coordinate chart this follows from Fourier inversion and Cauchy--Schwarz with the integrable weight $(1+|\xi|^2)^{-2}$, and a finite partition of unity gives the global bound. Consequently
\[
 \|v_n-w_n\|_\infty=O(\sqrt n).
\]
At the first boundary point with $s=\pi/n$ this yields
\[
 h(s)=h_0(s)-\tfrac12(n+n^{-1})+O(\sqrt n)<0
\]
for every sufficiently large $n$. Fix any such $n$. Since $a_1\ge1/2$, all three collars are everywhere inward and are embedded and disjoint for small $\varepsilon$. Their boundaries are analytic. The variation formula then gives $u_\varepsilon<0$ at the corresponding boundary point for all small positive $\varepsilon$.

Finally the old area of the retained surface is strictly less than $-2\pi\chi(M)$, while the new area equals that number. Hence $\int_{M'}(e^{2u_\varepsilon}-1)dA_g>0$, which forces $u_\varepsilon>0$ somewhere. This proves that no pointwise order holds in this family, and completes the general-shape answer to 67/021.
\end{proof}

\subsection{A completely specified collar counterexample}

\begin{proposition}[Proved: fixed frequency and fixed collar width]\label{A2-collar-explicit}
Let $M$ be the hyperbolic pant with all three geodesic boundary lengths equal to $2\pi$. Label the boundaries $B_1,B_2,B_3$, and take an arclength coordinate $s\in\R/(2\pi\Z)$ on $B_1$, with origin at its shortest perpendicular to $B_2$. Set
\[
 n=10^4,\qquad \varepsilon=10^{-14},\qquad
 a_1(s)=1+\tfrac12\cos(ns),\qquad a_2=a_3=1.
\]
In inward Fermi coordinates remove the three collars $0\le r<\varepsilon a_i(s)$. They are disjoint, everywhere inward, and have analytic boundaries. On the retained pant let $g'=e^{2u}g$ be its conformally compatible hyperbolic metric with geodesic boundary. At the point of its first boundary with $s=\pi/n$ one has
\[
 u\le-\frac{9991}{20}\,10^{-14}<0.
\]
The same function $u$ is positive somewhere. Thus the general-shape counterexample of Proposition~\ref{A2-collar-shape} can be given with these fixed parameters, without an unspecified perturbation threshold.
\end{proposition}

\begin{proof}
Put $y=\arctan(\sinh r)$, so
\[
 g=\sec^2y\,(ds^2+dy^2),\qquad
 \Delta_g=\cos^2y\,(\partial_s^2+\partial_y^2),\qquad
 L=\Delta_g-2.
\]
The collar theorem, applied to the double of $M$, gives disjoint collars of Fermi width $\operatorname{arsinh}(1/\sinh\pi)$; see \cite[Theorem~2.18, pp.~22--23]{McMullen1994A567}. In particular the three closed collars $0\le y\le1/50$ are embedded and disjoint. Indeed $\pi<4$, $e<3$, and
\[
 \tan(1/50)\le\frac{1/50}{1-(1/50)^2/2}
 =\frac{100}{4999}<\frac2{81}<\frac1{\sinh\pi}.
\]
The selected new boundaries lie in these collars, and even in $y<1/100$, since $0<\varepsilon a_i\le3\varepsilon/2<1/100$ and $y\le r$.

Here are explicit comparison functions. Write $b=1/100$, and define the $C^2$ cutoff
\[
 \eta(y)=\begin{cases}
 1,&0\le y\le b,\\
 1-10t^3+15t^4-6t^5,&b\le y\le2b,\quad t=(y-b)/b,\\
 0,&y\ge2b.
 \end{cases}
\]
Its first two derivatives vanish at the endpoints of the transition interval. On that interval,
\[
 |\eta'|\le\frac{15}{8b}<200,\qquad
 |\eta''|\le\frac{10}{\sqrt3\,b^2}<60000.
\]
These bounds follow by differentiating the displayed polynomial: its derivatives are $-30t^2(1-t)^2$ and $-60t(1-t)(1-2t)$. Their exact maximum absolute values on $[0,1]$ are $15/8$ and $10/\sqrt3$. Also, throughout $0\le y\le2b$,
\[
 \tan y<\frac1{40},\qquad \sec^2y<2.
\]
For the second inequality use $\cos y\ge1-y^2/2\ge4999/5000$.

Extend each $\eta(y_i)\tan y_i$ by zero outside its collar and set
\[
 H=2000-\sum_{i=1}^3\eta(y_i)\tan y_i.
\]
At most one summand is nonzero at a point. Since $L(\tan y)=0$, differentiation gives
\[
 LH=-4000-\cos^2y\bigl(\eta''\tan y+2\eta'\sec^2y\bigr)
 \le-4000+60000/40+800=-1700
\]
on a collar, and $LH=-4000$ elsewhere.

The separated function
\[
 f_n(y)=e^{-ny}(n-\tan y)
\]
satisfies the exact identity $f_n''-(n^2+2\sec^2y)f_n=0$ and $-f_n'(0)=n^2+1$. Define
\[
 w=\eta(y_1)f_n(y_1)\cos(ns)
\]
on the first collar and extend it by zero. Its residual is supported in $b\le y_1\le2b$ and equals
\[
 Lw=\cos^2y_1\bigl(\eta''f_n+2\eta'f_n'\bigr)\cos(ns).
\]
Consequently, using $nb=100$ and $e>2$,
\begin{align*}
 \|Lw\|_\infty
 &\le e^{-100}\left[60000(n+1/40)
       +400(n^2+n/40+2)\right]\\
 &<\frac{40600102300}{2^{100}}<1.
\end{align*}
All bounds so far are explicit rational or integer inequalities apart from the elementary displayed trigonometric estimates.

Set $\Psi=2H+w/2$ and $U=\varepsilon\Psi$. Then $L\Psi\le-3400+1/2<0$. Since $e^{2x}-1\ge2x$ for every real $x$, this proves
\[
 \Delta_g U\le e^{2U}-1.
\]
It remains to check the outward normal derivative on each new boundary; an interior supersolution alone would not suffice.

Write a boundary graph as $r=f(s)=\varepsilon a(s)$, and put $p=f'(s)$, $c=\cosh f$, $q=p/c$. Its outward normal and signed geodesic curvature are exactly
\[
 N=\frac{-c\,\partial_r+(p/c)\partial_s}{\sqrt{c^2+p^2}},\qquad
 k_g=\frac{\tanh f(1+2q^2)-f''\operatorname{sech}^2f}
                {(1+q^2)^{3/2}}.
\]
For example, the curvature formula follows by inserting the graph velocity $f'\partial_r+\partial_s$ into $g(\nabla_TT,N)$, using $\Gamma^r_{ss}=-\cosh r\sinh r$ and $\Gamma^s_{rs}=\tanh r$.

We give the error bounds used in the comparison. On the region $0\le r\le1/100$, the cutoff is one. On the first collar there,
\[
 \Psi=4000-2\sinh r+\tfrac12F_n(r)\cos(ns),\qquad
 F_n(r)=e^{-n\arctan(\sinh r)}(n-\sinh r).
\]
On the other two collars the oscillating term is absent. Because $\sinh r<1$, $\cosh r<2$, $|y'|\le1$, and $|y''|\le1$, direct differentiation gives
\[
 |F_n|\le2n,\quad |F_n'|\le2n^2,\quad
 |F_n''|\le n^3+2n^2+5n+1\le4n^3.
\]
It follows that
\[
 |\Psi_r|\le2n^2,\qquad |\Psi_s|\le n^2,
 \qquad |\Psi_{rr}|\le3n^3.
\]
At $r=0$ the outward derivative is
\[
 -\Psi_r(0,s)=a(s)-a''(s)+1.
\]
The inequalities $f\le3\varepsilon/2$, $|p|\le\varepsilon n/2$, and
$1-(1+q^2)^{-1/2}\le q^2/2$ therefore imply
\begin{align*}
 |\partial_N\Psi-(a-a''+1)|
 &\le 3n^3f+n^2p^2+n^2|p|\\
 &\le5\varepsilon n^3+\tfrac14\varepsilon^2n^4
 <10\varepsilon n^3.
\end{align*}
The same estimates apply on the other two boundary components.

For the curvature, use $|\tanh f-f|\le f^3/3$,
$1-\operatorname{sech}^2f\le f^2$, and
$1-(1+x)^{-3/2}\le3x/2$ for $x\ge0$. On $0\le x\le1$ the derivative of $(1+2x)(1+x)^{-3/2}$ has absolute value at most $1/2$, so its distance from one is at most $2x$. Since $q^2\le p^2<1$, the exact curvature formula yields
\begin{align*}
 |k_g-\varepsilon(a-a'')|
 &\le f^3/3+2fq^2+|f''|(f^2+3q^2/2)\\
 &\le\varepsilon^3\left(\tfrac98+\tfrac{15}8n^2+\tfrac3{16}n^4\right)
 \le\varepsilon^3n^4.
\end{align*}
Thus, for the specified $n,\varepsilon$,
\[
 \partial_NU-k_g
 \ge\varepsilon\bigl(1-10\varepsilon n^3-\varepsilon^2n^4\bigr)>0,
 \qquad 10\varepsilon n^3+\varepsilon^2n^4=\frac1{10}+10^{-12}<1.
\]

The function $u$ satisfies $\Delta_gu=e^{2u}-1$ and $\partial_Nu=k_g$. Apply the maximum principle and boundary-point lemma to $U-u$: it satisfies $(\Delta_g-c_0)(U-u)\le0$ with $c_0>0$ and outward derivative nonnegative. A negative minimum is impossible, whether interior or on the boundary. Hence $u\le U$ on the retained pant.

At the specified boundary point, $f=\varepsilon/2$ and $\cos(ns)=-1$. Here $ny\le n\varepsilon/2<1/10$, so $e^{-ny}\ge1-ny\ge9/10$, while $\sinh f<1$. Therefore
\[
 \Psi\le4000-\tfrac9{20}(n-1)=-\frac{9991}{20}.
\]
This proves the stated strict negative bound. The new boundaries are analytic because their defining functions $a_i$ are analytic; the merely $C^2$ cutoff occurs only in an auxiliary comparison function. Their positive inward widths and the explicit disjoint collars verify every geometric hypothesis. Finally the retained surface has old area less than $2\pi$, while its new geodesic-boundary hyperbolic area is $2\pi$ by Gauss--Bonnet. Hence $u>0$ somewhere. The exact symbolic identities and rational inequalities used above are checked in \path{code/part_a567_collar_explicit.py}; no numerical approximation to $u$ is used.
\end{proof}

